% ══════════════════════════════════════════════════════════════════════════
\section{Introduction}
\label{sec:intro}
% ══════════════════════════════════════════════════════════════════════════

The standing counsel to the Bitcoin holder who fears physical coercion has three
parts: hold quietly, deny ownership if asked, and keep something to surrender if
denial fails. Each is advice about what an attacker will \emph{believe}. None of
it has been assessed as what it is --- a security mechanism whose strength
depends entirely on the adversary's ignorance, which is the weakness
\emph{Reliance on Security Through Obscurity} catalogs~\cite{cwe656}.

This paper assesses it, and reaches a conclusion stronger than ineffectiveness.
The counsel is \emph{self-defeating in aggregate}, by a mechanism with the
structure of a commons. Denial credibility is supplied by the population of
holders and consumed by each holder who denies. Where denial becomes the expected
script, a denial stops carrying information; an interrogation that opens with one
yields the attacker no exculpatory update; and its rational continuation is
further coercion rather than release. Each holder's individually rational
response to that environment --- conceal harder, rehearse better --- depletes the
resource further, for everyone. We call this the \textbf{Denial Spiral}
(\S\ref{sec:spiral}), and its cost is an externality: it falls on whoever must be
believed, including the holder of a design that depletes nothing, and including
the person who holds no bitcoin at all.

\paragraph{Relation to the companion paper.} \cite{deadly-race} derives a bar for
coercion-resistant custody --- the \emph{No-Gray-Area Criterion} --- from the
observation that a seizure leaving the attacker a \emph{feasible-but-unfinished}
path to the funds makes the holder a live competing racer, and so prices his
removal. That is the \emph{Deadly Race}: it governs the \emph{terminal} move. The
Spiral governs the \emph{duration}. The two are companions rather than one being
a case of the other: a design can be safe under one and guilty under the other,
and the wallet-erase PIN of \S\ref{sec:obscurity} is exactly such a design. This
paper is self-contained --- \S\ref{sec:model} restates what it borrows --- but
the formal development of the adversary, the criterion and the construction that
clears it are in the companion and are not repeated.

\paragraph{Contributions.}
\begin{enumerate}[nosep,leftmargin=1.9em]
  \item The \textbf{Denial Spiral}: obscurity counsel as a depleting commons,
        with its cost falling outside the population that depletes it
        (\S\ref{sec:spiral}).
  \item A \textbf{classification of the shipped obscurity class} --- wallet-erase
        and duress PINs, decoy passphrase wallets, and concealment marketed as
        deniability --- against the Kerckhoffs premise, on vendors' own
        documentation (\S\ref{sec:obscurity}).
  \item An account of \textbf{what kind of claim this is}: derived from purposive
        action rather than induced from observation, with an argument that no
        observation could settle it, because every witness to an enacted case has
        an interest in the answer (\S\ref{sec:method}).
  \item \textbf{Four censoring mechanisms} that keep the crisis undiagnosable,
        the fourth of which removes the response to the evidence and so survives
        the evidence being published (\S\ref{sec:anosognosia}).
  \item \textbf{Registry coding} of $352$ incidents that grounds the stipulations
        and bounds what the argument may claim (\S\ref{sec:registry}).
\end{enumerate}
